\lectures{14}
\title{NP-Completeness}

\begin{document}

\subsection{NP-Completeness}

\textbf{Definition.}
A language $B$ is \textbf{NP-complete}
if $B \in \mathrm{NP}$ and
$A \leq_p B$ for every $A \in \mathrm{NP}$.

\begin{remark}
    Make an analogy:
    the language $A_{\mathrm{TM}}$ is Turing-recognizable,
    and $L \leq_m A_{\mathrm{TM}}$ for any Turing-recognizable $L$.
\end{remark}

The star of NP-completeness will be $\mathrm{3SAT}$.

For today's lecture,
let's black-box the following result;
we'll prove this result next lecture.

\begin{boxed-text}
    \textbf{Cook-Levin Theorem.}
    The language $\mathrm{3SAT}$ is
    NP-complete.
\end{boxed-text}

\begin{remark}
    Refer to \href{/6.1220/lectures/14/}{6.1220 Notes}
    for a definition of $\mathrm{3SAT}$.
\end{remark}

We'll use the above result to prove the NP-completeness
of other languages.

\subsection{CLIQUE}

\begin{theorem}{CLIQUE}
    The language $\mathrm{CLIQUE} := \{
        \langle G, k \rangle \mid
        G \text{ is an undirected graph with a } k\text{-clique}
    \}$ is NP-complete.
\end{theorem}

\begin{proof}
    It suffices to construct a
    poly-time reduction from $\mathrm{3SAT}$ to $\mathrm{CLIQUE}$.
    Given a 3CNF with $k$ clauses, we construct $\langle G, k \rangle$, where $G$ is a graph in which:
    \begin{itemize}
        \item The nodes are the literals, one per occurrence in a clause.
        \item An edge is drawn between two literals if and only if
        they are \emph{not} \dots
        \begin{itemize}
            \item \dots in the same clause, nor \dots
            \item \dots complements of each other
            (e.g., $a$ and $\neg a$).
        \end{itemize}
    \end{itemize}
    You can play around with the (warning: AI-generated) widget below
    to get a feel for why this reduction works.

    \includewidget{lecture-14/clique}

    Note that one must check two things to show the reduction is valid:
    both $w \in A \implies f(w) \in B$ and $f(w) \in B \implies w \in A$.
    Convince yourself that both of these are true. ~ $\blacksquare$
\end{proof}

\subsection{HAMPATH}

\begin{theorem}{HAMPATH}
    The language $\mathrm{HAMPATH} := \{
        \langle G, s, t \rangle \mid
        G \text{ is a directed graph with a Hamiltonian path from } s \text{ to } t
    \}$ is NP-complete.
\end{theorem}

\begin{proof}
    Again, we construct a poly-time reduction
    from $\mathrm{3SAT}$ to $\mathrm{HAMPATH}$.
    Given a 3CNF, we construct a graph $G$ in which:
    \begin{itemize}
        \item There is one \emph{diamond} for every variable $x_i$.
        \begin{itemize}
            \item If $x_i = \textsc{True}$,
            then the diamond's middle path points from left to right.
            If $x_i = \textsc{False}$, it's right to left.
        \end{itemize}
        \item There is one node for every clause.
        \begin{itemize}
            \item \emph{Example:} If clause $C_2$ has literals
            $\neg x_1 \vee \neg x_2 \vee x_3$,
            then it may only be reached by \dots
            \begin{itemize}
                \item \dots traversing $x_1$'s diamond
                from right to left (i.e., if $x_1 = \textsc{False}$).
                \item \dots traversing $x_2$'s diamond
                from right to left (i.e., if $x_2 = \textsc{False}$).
                \item \dots traversing $x_3$'s diamond
                from left to right (i.e., if $x_3 = \textsc{True}$).
            \end{itemize}
        \end{itemize}
    \end{itemize}
    You can play around with the (warning: AI-generated) widget below
    to see how this works.

    \includewidget{lecture-14/hampath}

    As before, a formal proof would justify this reduction
    in both directions, but we omit the details. ~ $\blacksquare$
\end{proof}

\end{document}